% kaobook starter — a wide-margin book/notes class with sidenotes and a
% mini table of contents at the start of each chapter.
%
% kaobook.cls, kao.sty, kaorefs.sty and kaotheorems.sty are VENDORED here
% (LPPL-1.3c, (c) Francesco Marotta) because kaobook is not a CTAN package —
% see the README for why and how to update them.
\documentclass[
	fontsize=10pt, % Base font size
	twoside=false, % Use different layouts for even and odd pages
	secnumdepth=1, % How deep to number headings
]{kaobook}

\usepackage[english]{babel}
\usepackage[english=british]{csquotes}

% Cross-references and margin notes, in that load order (kaorefs must load
% after kaotheorems if both are used, and before hyperref is configured).
\usepackage{kaotheorems}
\usepackage{kaorefs}

\graphicspath{{images/}{./}}

\makeindex[columns=2, title=Index, intoc]

\begin{document}

\titlehead{Lecture Notes}
\title[Notes on the Subject]{Notes on the Subject}
\author[J. Doe]{Jane Doe}
\date{\today}
\publishers{Department of Whatever}

\frontmatter

\makeatletter
\uppertitleback{\@titlehead}
\lowertitleback{
	\textbf{About this document} \\
	Replace this page with your own front matter --- a colophon, a licence
	statement, printing history, or nothing at all.

	\medskip

	\textbf{Colophon} \\
	Typeset with \href{https://sourceforge.net/projects/koma-script/}{\KOMAScript}
	and \href{https://www.latex-project.org/}{\LaTeX} using the
	\href{https://github.com/fmarotta/kaobook/}{kaobook} class.
}
\makeatother

\maketitle

\chapter*{Preface}
Say what this document is, who it is for, and how it is organised. This
paragraph is the first thing a reader sees after the title page, so make it
count.

\begingroup
\etocstandarddisplaystyle
\etocstandardlines
\tableofcontents
\endgroup

\mainmatter
\setchapterstyle{kao}

\chapter{Getting Started}
\label{ch:intro}
Start the first chapter here. kaobook keeps a mini table of contents in the
margin for every chapter, updated automatically from your section headings.

\section{Sidenotes}
Use a sidenote for anything that would break the flow of the main
text\sidenote{Like this. It appears in the wide margin, roughly level with
where it is called.} but is still worth keeping close to the point it
supports.

\section{Cross-references}
Refer back to earlier material with \cref{ch:intro}, or forward to
\cref{ch:next} once it exists --- \texttt{kaorefs} loads \texttt{cleveref}
configured for kaobook's numbering.

\section{Figures and tables}
Figures and tables belong in the margin by default in this class --- use
\texttt{marginfigure} and \texttt{margintable} rather than the plain
\texttt{figure}/\texttt{table} environments, which is also where kaobook's
own caption styling is set up to work.\sidenote{Plain \texttt{figure} and
\texttt{table} still typeset, but the class does not style their captions the
way it styles a margin float's --- use \texttt{pagelayout\{wide\}} instead if
you want a full-width figure.}

\begin{marginfigure}
	\centering
	\fbox{\parbox{0.9\linewidth}{\centering\vspace{1.2cm}Your figure here\vspace{1.2cm}}}
	\caption{A caption goes below the figure.}
	\label{fig:example}
\end{marginfigure}

\cref{fig:example} shows a placeholder. Replace the \texttt{fbox} with
\texttt{\textbackslash includegraphics}.

\begin{margintable}
	\centering
	\begin{tabular}{lr}
		\toprule
		Condition & Rate \\
		\midrule
		Baseline & 0.42 \\
		Ours     & 0.71 \\
		\bottomrule
	\end{tabular}
	\caption{A small margin table.}
	\label{tab:example}
\end{margintable}

\cref{tab:example} shows the same idea for a table.

\chapter{Building on It}
\label{ch:next}
A second chapter, so the mini table of contents and cross-references above
have somewhere real to point.

\pagelayout{wide}
\addpart{Appendix}
\pagelayout{margin}

\appendix
\chapter{Supplementary Material}
Appendices are numbered with letters by convention, switched on by
\texttt{\textbackslash appendix}.

\backmatter
\setchapterstyle{plain}

\printindex

\end{document}
